


\chapter{Resultados Esperados}\label{chap:resultados_esperados}


Do agente de conformidade de gênero espera-se que a extração de \textbf{moves} concorde com a anotação de referência em proporção suficiente para sustentar as camadas que dela dependem, com os erros concentrados em categorias identificáveis e não dispersos por todo o esquema. Espera-se, em seguida, que a Especificação de Gênero derivada do corpus descreva o gênero da trilha, e não regularidades acidentais do conjunto processado. Três resultados principais são esperados: que artigos aceitos e retidos antes da derivação satisfaçam a maior parte das regras de estatuto obrigatório; que a conformidade das trilhas de contraste seja inferior à dos artigos retidos, indicando que as regras capturam algo próprio da trilha-alvo e não propriedades genéricas de artigos publicados; e que as especificações derivadas da \textit{Main track} e da \textit{Undergraduate track} do ENIAC difiram entre si, o que sustentaria a trilha como unidade de análise adequada. Espera-se ainda que a análise de sensibilidade identifique um núcleo de regras cujo estatuto permanece estável ao longo da faixa de limiares, distinguindo-o das regras que oscilam.

Do agente de fronteira de rejeição espera-se que os motivos de recusa extraídos das avaliações correspondam aos que os revisores efetivamente apontaram, medida pela concordância com a anotação de referência, e que o retorno de risco produzido sobre a partição retida do ReviewArena distinga artigos recusados de artigos aceitos com desempenho superior ao de uma linha de base trivial. Espera-se que esse resultado estabeleça que avaliações negativas reais permitem sinalizar risco de recusa para um manuscrito ainda não submetido. Cabe delimitar seu alcance: o gabarito provém de conferências internacionais de aprendizado de máquina e processamento de linguagem natural, de modo que um resultado favorável demonstraria o mecanismo, e não sua transferência a uma trilha brasileira.

Como resultado de alcance mais amplo, espera-se que o artefato seja replicável a outras trilhas, outras edições e outras conferências, uma vez que incluir um novo conjunto consiste em reexecutar os fluxos dos agentes, sem alteração no restante do artefato. Espera-se, pelo mesmo motivo, que a arquitetura comporte a substituição do ReviewArena por um corpus de avaliações brasileiras, quando este existir, e a incorporação de agentes adicionais.

No plano teórico, espera-se consolidar a caracterização das convenções de escrita de uma comunidade científica como um gênero operacionalizável por um artefato computacional.

Um resultado adverso em qualquer dos níveis é achado da pesquisa, e não falha a ser contornada. Estabelecer que a modelagem de gênero por derivação empírica não distingue trilhas de uma mesma comunidade, ou que avaliações negativas de outras comunidades não permitem sinalizar risco de recusa, seria contribuição legítima ao conhecimento, ainda que negativa, e indicaria com precisão qual camada do artefato revisar.

% Do agente de conformidade de gênero espera-se, que a extração de \textbf{moves} concorde com a anotação de referência em proporção suficiente para sustentar as camadas que dela dependem, com os erros concentrados em categorias identificáveis e não dispersos por todo o esquema. Espera-se, em seguida, que a Especificação de Gênero derivada do corpus descreva o gênero da trilha, e não regularidades acidentais do conjunto processado. Três resultados principais são esperados: que artigos aceitos e retidos antes da derivação satisfaçam a maior parte das regras de estatuto obrigatório; que a conformidade das trilhas de contraste seja inferior à dos artigos retidos, indicando que as regras capturam algo próprio da trilha-alvo e não propriedades genéricas de artigos publicados; Que o agente esteja preparado para consolidar e construir as regras de generos ed outras trilhas do SBC sem a necessidade alteração do escopo. 

% Do agente de fronteira de rejeição espera-se que os motivos de recusa extraídos das avaliações correspondam aos que os revisores efetivamente apontaram, medida pela concordância com a anotação de referência, e que o retorno de risco produzido sobre a partição retida do ReviewArena distinga artigos recusados de artigos aceitos com desempenho aceitável com a proposta do artefato. Espera-se que ele estabeleça que avaliações negativas reais permitem sinalizar risco de recusa para um manuscrito ainda não submetido. Espera-se que a arquitetura proposta esteja preparada para comprotar a troca de corpus e a incorporação de agentes adicionais sem revisão da arquitetura, e que o artefato seja capaz de operar sobre outras trilhas do SBC, com a substituição do ReviewArena por um corpus de avaliações brasileiras quando este existir.

% Como resultado de alcance mais amplo, espera-se que o modelo proposto como projeto final seja replicável a outras trilhas, outras edições e outras conferências, uma vez que incluir um novo conjunto consiste em reexecutar os fluxos dos agentes, sem alteração no restante do artefato. 

% No plano teórico, espera-se consolidar a caracterização das convenções de escrita de uma comunidade científica como um gênero operacionalizável por um artefato computacional.

% Um resultado adverso em qualquer dos níveis é achado da pesquisa, e não falha a ser contornada. Estabelecer que a modelagem de gênero por derivação empírica não distingue trilhas de uma mesma comunidade, ou que avaliações negativas de outras comunidades não permitem sinalizar risco de recusa, seria contribuição legítima ao conhecimento, ainda que negativa, e indicaria com precisão qual camada do artefato revisar.